% ============================================================
%  Nota wyjaśniająca: Trzy prędkości światła w TGP
%  Status: Propozycja (do włączenia w sek04 jako rem:three-speeds)
% ============================================================

\begin{remark}[Trzy prędkości światła w TGP i ich relacje]%
\label{rem:three-speeds}

W~formaliźmie TGP współistnieją \textbf{trzy różne} wielkości
mające wymiar prędkości, które mogą być mylone z~,,prędkością
światła''. Poniższa tabela eliminuje tę niejednoznaczność.

\medskip
\begin{center}\small
\begin{tabular}{@{}lclc@{}}
\toprule
\textbf{Symbol} & \textbf{Definicja} & \textbf{Interpretacja}
  & \textbf{Wykładnik $a$} \\
\midrule
$c_{\mathrm{proper}}$ & $c_0$ (zawsze)
  & prędkość lokalna (zasada równoważności)
  & $a = 0$ \\[3pt]
$c_{\mathrm{lok}}(\Phi)$ & $c_0\sqrt{\Phi_0/\Phi}$
  & odl.~właściwa / czas koordynatowy (Aksjomat~A6)
  & $a = 1/2$ \\[3pt]
$c_{\mathrm{koord}}(\Phi)$ & $c_0\,\Phi_0/\Phi$
  & odl.~koordynatowa / czas koordynatowy (z metryki)
  & $a = 1$ \\
\bottomrule
\end{tabular}
\end{center}
\medskip

\noindent\textbf{Wyprowadzenie.}
Z~metryki efektywnej TGP:
\begin{equation}
  ds^2 = -\frac{c_0^2}{\psi}\,dt^2 + \psi\,\delta_{ij}\,dx^i dx^j,
  \qquad \psi = \Phi/\Phi_0,
\end{equation}
geodezyjka zerowa ($ds^2 = 0$) daje:
\begin{equation}
  \left(\frac{dr}{dt}\right)^2
  = \frac{c_0^2}{\psi^2}
  \quad\Rightarrow\quad
  c_{\mathrm{koord}} = \frac{c_0}{\psi} = c_0\,\frac{\Phi_0}{\Phi}.
\end{equation}
Prędkość mierzona jako \emph{odległość właściwa} (czynnik $\sqrt{h} = \sqrt{\psi}$)
podzielona przez \emph{czas koordynatowy}:
\begin{equation}
  c_{\mathrm{lok}} = \sqrt{h}\cdot\frac{dr}{dt}
  = \sqrt{\psi}\cdot\frac{c_0}{\psi}
  = \frac{c_0}{\sqrt{\psi}}
  = c_0\sqrt{\frac{\Phi_0}{\Phi}}.
\end{equation}
Prędkość mierzona w~\emph{obu właściwych miarach}
(właściwy czas: $d\tau = \sqrt{|g_{00}|/c_0^2}\,dt = dt/\sqrt{\psi}$):
\begin{equation}
  c_{\mathrm{proper}}
  = \frac{\sqrt{h}\,dr}{d\tau}
  = \frac{\sqrt{\psi}\cdot c_0/\psi}{1/\sqrt{\psi}}
  = c_0.
\end{equation}

\noindent\textbf{Relacja między trzema prędkościami:}
\begin{equation}\label{eq:three-c-relation}
  \boxed{
    c_{\mathrm{koord}} = \frac{c_{\mathrm{lok}}^2}{c_0}
    = \frac{c_0}{\psi},
    \qquad
    c_{\mathrm{proper}} = c_0
    \quad (\text{zawsze}).
  }
\end{equation}

\noindent\textbf{Spójność z~metryką:}
Aksjomat~A6 definiuje $c(\Phi) \equiv c_{\mathrm{lok}}$,
nie $c_{\mathrm{koord}}$.
Metryka daje $c_{\mathrm{koord}} = c_0/\psi$,
co jest \textbf{spójne}, ponieważ:
\begin{equation}
  g_{tt} = -\frac{c_0^2}{\psi} = -c_{\mathrm{lok}}^2,
  \qquad
  g_{ij} = \psi\,\delta_{ij},
  \qquad
  \frac{|g_{tt}|}{g_{rr}} = \frac{c_{\mathrm{lok}}^2}{\psi}
  = c_{\mathrm{koord}}^2.
\end{equation}
Zatem $c_{\mathrm{lok}} = \sqrt{|g_{tt}|}$ (do czynnika $c_0$),
a~$c_{\mathrm{koord}} = \sqrt{|g_{tt}|/g_{rr}}$. Oba wynikają
z~tej samej metryki, żadna niespójność nie występuje.

\medskip
\noindent\textbf{Konsekwencja:} W~testach obserwacyjnych
(Shapiro delay, ugięcie światła, lensing) obserwable zależą
od $c_{\mathrm{koord}}$ (pomiar czasu koordinatowego przy
odległości koordynatowej). To zgadza się z~PPN, gdyż
$c_{\mathrm{koord}} \approx c_0(1 - 2U/c_0^2)$
w~słabym polu.

\medskip
\noindent\textbf{Dlaczego A6 używa $c_{\mathrm{lok}}$, nie $c_{\mathrm{koord}}$?}
Ponieważ $c_{\mathrm{lok}}$ jest \textbf{jedyną
z~trzech prędkości, która jest (i) zmienna, (ii) niezależna
od wyboru układu koordynatowego}. Prędkość $c_{\mathrm{proper}} = c_0$
jest stała (nieinformatywna), a~$c_{\mathrm{koord}}$ zależy
od parametryzacji. Natomiast $c_{\mathrm{lok}}$ jest
skalarem: mierzy fizycznie, jak daleko światło zajedzie
w~jednostce czasu referencyjnego $dt$ --- to relacyjna
obserwabla TGP.
\end{remark}
